\section{The transfer and all seven final conclusions}\label{sec:conclusions}
\begin{theorem}[The second input]\label{thm:second}
For every full domain \(D\) on \(n\geq1\) alternatives,
\[
 \boxed{3nM\leq8(\Sig+G).}
\]
\end{theorem}
\begin{proof}
Apply Theorem~\ref{thm:density} to each component and sum. By
Theorem~\ref{thm:components}, the cardinalities sum to \(M\) and the
deletion sums to \(\Sig+G\). The empty-domain case has both sides zero.
\end{proof}
\lean{\path{CDCert.TwoInputs.second_input}, literally on
\(\mathrm{Fin}(k+1)\), with fullness as its only domain premise.}

\begin{theorem}[Transfer to general domains]\label{thm:transfer}
For every maximal Condorcet domain on \(n\geq1\) alternatives,
\[
 \boxed{(11n+8)M\leq32\Sig.}
\]
\end{theorem}
\begin{proof}
Maximality supplies fullness. The second input and the gap inequality
give \(3nM/8\leq\Sig+G\leq4\Sig-(n+1)M\). Multiply by eight
and rearrange: \((3n+8n+8)M\leq32\Sig\).
\end{proof}
\lean{\path{CDCert.TwoInputs.transfer} and
\path{CDCert.Recurrences.deletion_lower_bound}. There is no assumed
injection or density theorem in these statements: both have been
proved in the dependency chain. The displayed rational algebra is
performed over natural counts after clearing denominators.}

\begin{theorem}[The seven final entry-point conclusions]\label{thm:final}
Let \(f,h\) have the semantic definitions in Section~\ref{sec:definitions}.
The first four assertions hold for every integer \(n\geq3\).
\begin{enumerate}[label=E\arabic*.]
\item \((11n+8)f(n)\leq32n f(n-1)\).
\item \(f(n)\leq4(32/11)^{n-3}\).
\item \(3h(n)\leq8h(n-1)\).
\item \(h(n)\leq4(8/3)^{n-3}\).
\item There is a real number \(L_f\) such that
\(f(n)^{1/n}\to L_f\) and \(L_f\leq32/11\).
\item There is a real number \(L_h\) such that
\(h(n)^{1/n}\to L_h\) and \(L_h\leq8/3\).
\item \(f(n)=O\bigl(n^{-8/11}(32/11)^n\bigr)\) as \(n\to\infty\).
\end{enumerate}
\end{theorem}
\lean{These are exactly, in the same order,
\path{CDCert.Main.general_recurrence}, \path{general_geometric},
\path{peak_recurrence}, \path{peak_geometric}, \path{general_limit},
\path{peak_limit}, \path{general_polynomial_correction}.
The limit assertions quantify existentially over a real limit; they
assert neither an exact value nor convergence of successive ratios.
The final three statements have no extra mathematical hypotheses.}

\subsection{Recurrences and the base value}
\begin{proof}[Proof of E1--E4]
Choose a domain of size \(f(n)\). It is maximal, so
Theorem~\ref{thm:transfer} and Lemma~\ref{lem:deletion-upper} give E1.
For a maximum peak-pit domain, completion puts it in \(\F_n\).
The density theorem and the peak-pit deletion estimate give E3.

On three alternatives a common never-law excludes two of the six
orders, so every Condorcet domain has at most four orders. The
peak-pit example in Section~\ref{sec:deletion} attains four. Hence
\(f(3)=h(3)=4\). Starting there, the general recurrence factor satisfies
\(32n/(11n+8)\leq32/11\), so induction gives E2. Induction using
E3 gives E4. The shift by three and multiplier four come from this
base value. The shifted formulas are not claimed at \(n=1,2\), where
both functions take the direct values 1 and 2.
\end{proof}
\lean{\path{Condorcet.f_three}, \path{h_three};
\path{CDCert.Recurrences.general_recurrence'} and
\path{CDCert.Recurrences.peakPit_recurrence} actually give the recurrences for
\(n\geq1\). The final \path{Main} wrappers use the paper's \(n\geq3\)
interface. \path{f_geometric}, \path{CDCert.Recurrences.h_geometric} first prove the
integer-scaled power bounds; \path{CDCert.Asymptotics.f_le_real} and
\path{h_le_real} give the real-valued forms.}

\subsection{Root-growth limits: a separate use of concatenation}
\begin{proof}[Proof of E5 and E6]
A singleton domain shows \(f(n),h(n)\geq1\). Ordered concatenation of
any two Condorcet domains is again Condorcet: majority restricts to
transitive majorities inside the blocks and always puts the first
block above the second. For peak-pit domains, a mixed triple has a
label fixed first or last, so an extreme law exists. Cardinalities
multiply. Thus, for \(g=f\) or \(h\),
\[
 g(m+n)\geq g(m)g(n).
\]
Put \(a_n=\log g(n)\geq0\); it is superadditive. For a fixed
\(k\geq1\), write \(n=qk+r\), \(0\leq r<k\). Repeated
superadditivity gives \(a_n\geq q a_k+a_r\geq q a_k\), hence
\(\liminf a_n/n\geq a_k/k\). Taking the supremum over \(k\)
shows that the liminf is at least the supremum of all \(a_k/k\),
whereas the limsup cannot exceed that supremum. The exponential
upper bound makes it finite. Therefore \(a_n/n\) has a limit.
The bounds E2 or E4 imply that this limit is at most \(\log(32/11)\)
or \(\log(8/3)\). Exponentiating proves E5 and E6.
\end{proof}
\lean{\path{CDCert.Supermult.f_supermul}, \path{h_supermul},
\path{one_le_f}, \path{one_le_h}; then
\path{CDCert.Asymptotics.root_limit}, \path{f_root_limit},
\path{h_root_limit}. The generic root-limit proof uses Mathlib's
subadditive machinery on the negative logarithm. This limit argument
is not needed to obtain E1--E4.}

\subsection{The polynomial improvement, with a direct proof}
\begin{proof}[Proof of E7]
Set \(r=32/11\), \(a=8/11\). The recurrence is
\[
 f(n+1)\leq r\frac{n+1}{n+1+a}f(n).
\]
For \(0<a<1\) and \(x\geq0\), the elementary concavity inequality
\((1+x)^a\leq1+ax\) follows by differentiating
\(1+ax-(1+x)^a\): its derivative is nonnegative and its value at zero
is zero. With \(x=1/(n+1)\), it gives
\[
 (n+2)^a(n+1)\leq(n+1)^a(n+1+a).
\]
Consequently \(f(n)(n+1)^a/r^n\) is non-increasing for \(n\geq3\).
Its initial value is \(C=4\cdot4^a/r^3\), so
\(f(n)\leq C(n+1)^{-a}r^n\leq Cn^{-a}r^n\) for all \(n\geq3\).
This is precisely the asserted big-O estimate.
\end{proof}
\lean{\path{CDCert.Asymptotics.bernoulli_step}, \path{cW},
\path{f_weighted}, \path{f_isBigO}; the last is restated by
\path{CDCert.Main.general_polynomial_correction}.
The real-power exponent is exactly \(-8/11\), not a natural-number
subtraction. The finitely many initial values are irrelevant to
\path{Asymptotics.IsBigO} at \texttt{atTop}.}
